\begin{titlepage}
\centering
{\Large Notas didácticas de una entrevista en profundidad\par}
\vspace{1.0cm}
{\huge\bfseries Diez años de evolución de la conducción autónoma: de los mapas de alta precisión y BEV a los modelos de extremo a extremo y los modelos del mundo}\par
\vspace{0.5cm}
{\Large Zhang Xiaojun Podcast: conversación con Lang Xianpeng (郎咸朋), de Li Auto (理想汽车)\par}
\vspace{0.35cm}
{\large Elaborado a partir del video público, la estructura de capítulos, la transcripción, la descripción del video y diagramas conceptuales propios\par}
\vspace{0.3cm}
{\large 2025-03-16, entrevista grabada en 2024-12\par}
\vspace{0.85cm}
\includegraphics[width=0.88\textwidth,height=0.42\textheight,keepaspectratio]{cover.jpg}\par
\vfill
\begin{tcolorbox}[width=0.92\textwidth, colback=black!2!white, colframe=black!55, sharp corners]
\textbf{Título del video}: 96. Conversación con Lang Xianpeng: diez años de evolución de la conducción autónoma, detalles técnicos clave y Tesla\par
\textbf{Canal del video}: Zhang Xiaojun Podcast\par
\textbf{Duración del video}: 02:00:20\par
\textbf{Enlace del video}: \href{\videourl}{\nolinkurl{\videourl}}\par
\textbf{Nota sobre los materiales}: YouTube no ofrece subtítulos descargables; el texto principal se elaboró a partir de una transcripción por bloques con faster-whisper large-v3 (CPU, int8), la descripción del video, la estructura de capítulos y figuras didácticas propias. Las tomas fijas de la entrevista no se reutilizan en el texto principal.\par
\textbf{Página del proyecto}: \href{\repourl}{\nolinkurl{\repourl}}\par
\end{tcolorbox}
\end{titlepage}

\begingroup
\small
\setlength{\parskip}{0pt}
\renewcommand{\baselinestretch}{0.92}\selectfont
\tableofcontents
\endgroup
\newpage
